\documentclass[12pt]{article}
\usepackage{amsmath,amssymb,amsthm}
\usepackage{hyperref}
\usepackage{booktabs}

\newtheorem{theorem}{Theorem}[section]
\newtheorem{lemma}[theorem]{Lemma}
\newtheorem{corollary}[theorem]{Corollary}
\newtheorem{proposition}[theorem]{Proposition}
\newtheorem{definition}[theorem]{Definition}
\newtheorem{remark}[theorem]{Remark}

\title{The $W(3,3)$ Graph: Uniqueness, Tau Values,\\
       and the CM $j$-Tower}
\author{Wil Dahn}
\date{\today}

\begin{document}
\maketitle

\begin{abstract}
The Ramanujan bipartite graph $W(3,3)$ of Lubotzky--Phillips--Sarnak is
characterized among the family $W(3,q)$ by six simultaneous algebraic
coincidences (C1--C6), each proved symbolically.
The spectral parameters encode the Ramanujan tau function at the first two primes:
$\tau(2)=-f(3)=-24$ and $\tau(3)=k(3)\Phi_6(3)\cdot 3=252$.
The $j$-invariants of all CM orders attached to the Ihara zeta sectors of $W(3,3)$
are perfect cubes of graph parameters, forming a five-entry tower
$\{k^3,\,-g^3,\,(v/2)^3,\,-\Sigma^3,\,255^3\}$ with $\Sigma=2^{\Phi_4/2}$;
this tower is generated entirely by the $s$-eigenspace sector, whose CM discriminant
$-4\Phi_6=-28$ is Heegner while the $r$-eigenspace discriminant $-4\Phi_4=-40$
has class number $2$. The sixth coincidence C6 --- that $3\mid g(3)$ uniquely
among prime $q$ --- makes $\Sigma=2^{g/3}$ an integer and links the $D=-11$
$j$-tower entry to the $r$-sector parameter $\Phi_4$ via $g/3=\Phi_4/2$.
As a physics application, the $W(3,3)$ spectral descent map predicts
a quasi-degenerate neutrino mass entry point $\mu_{\rm eff}^2=1/4$ and
seesaw cascade ratio $r=0.52244$. All results are verified by
$\mathbf{456}$ automated unit tests~\cite{Dahn2026repo}.
\end{abstract}

\tableofcontents
\bigskip

%%% INTRODUCTION %%%
\section*{Introduction}
\addcontentsline{toc}{section}{Introduction}
\label{sec:intro}

The Lubotzky--Phillips--Sarnak graph $W(3,3)$ is the unique member of the
infinite family $\{W(3,q)\}_{q\equiv 1(4)}$ for which the following six
algebraic coincidences hold simultaneously:
\begin{itemize}
\item[\textbf{C1.}] The $r$-eigenvalue discriminant equals $-4\Phi_4$:
  the eigenvalue equation $\mathrm{ev}_r^2-4(k-1)=-4\Phi_4$ holds only at $q=3$.
\item[\textbf{C2.}] $k+g = q^q=27$: degree plus $s$-multiplicity equals $3^3$.
\item[\textbf{C3.}] $-f=\tau(2)=-24$ and $kq\Phi_6=\tau(3)=252$:
  the spectral multiplicities encode the Ramanujan tau function at the first two primes.
\item[\textbf{C4.}] $\Phi_6(3)=7$ is a Heegner number:
  the CM field $\mathbb{Q}(\sqrt{-7})$ of the $s$-sector has class number $1$.
\item[\textbf{C5.}] $a_{k-1}(E)=\mathrm{ev}_s$: the Frobenius trace at $p=k-1=11$
  on the CM curve $E\colon y^2=x^3-35x+98$ ($j=-g^3$) equals the $s$-eigenvalue.
\item[\textbf{C6.}] $3\mid g(3)$: among all prime $q$, only $q=3$ has $3\mid g(q)$,
  making $\Sigma=2^{g/3}=2^{\Phi_4/2}$ an integer spectral invariant.
\end{itemize}
Together C1--C6 characterize $W(3,3)$ uniquely in $\{W(3,q)\}$
(Theorems~\ref{thm:uniqueness}--\ref{rem:C6}).

A consequence of C4--C6 is a five-entry \emph{$j$-tower}: the $j$-invariants
of all CM orders attached to the Ihara zeta sectors are perfect cubes of graph
parameters, $\{k^3,-g^3,(v/2)^3,-\Sigma^3,P^3\}$ with $P=255$
(Proposition~\ref{prop:cube}). No Heegner discriminant with $|D|>4\Phi_6=28$
contributes a cube root to the $W(3,3)$ spectral ring.

As a physics application (Sections~\ref{sec:neutrino}--\ref{sec:seesaw}),
the spectral descent map of $W(3,3)$ predicts a quasi-degenerate neutrino mass
entry point $\mu_{\rm eff}^2=1/4$ with seesaw cascade ratio $r=0.52244$,
in tension with the DESI DR1 bound $\sum m_\nu<72\,\mathrm{meV}$~\cite{DESI2024}.

All results are verified by $\mathbf{456}$ automated unit tests~\cite{Dahn2026repo}.

\medskip
\noindent\textit{Structure.}
\S\ref{sec:graph} defines $W(3,3)$ and its parameter table.
\S\ref{sec:uniqueness} proves C1--C5 symbolically.
\S\ref{sec:tau} establishes C3, C5--C6, the tau bridge, and the $j$-tower.
\S\S\ref{sec:neutrino}--\ref{sec:seesaw} develop the neutrino application.

%%% SECTION 1 %%%
\section{The Graph $W(3,3)$}
\label{sec:graph}

The Lubotzky--Phillips--Sarnak graph $W(3,q)$ \cite{LPS1988} is a
$k$-regular bipartite Ramanujan graph defined for prime powers
$q\equiv 1\pmod{4}$, with $k=q(q+1)$ and $|V|=2(q+1)(q^2+1)=2v$.
We work at $q=3$.

\begin{table}[h]
\centering
\begin{tabular}{lll}
\toprule
Parameter & Symbol & Value \\
\midrule
Degree              & $k=q(q+1)$          & $12$ \\
Vertices per side   & $v=(q+1)(q^2+1)$    & $40$ \\
$r$-eigenvalue      & $r=q-1$             & $2$ \\
$s$-eigenvalue      & $s=-(q+1)$          & $-4$ \\
$r$-multiplicity    & $f=q(q+1)^2/2$      & $24$ \\
$s$-multiplicity    & $g=q(q^2+1)/2$      & $15$ \\
$\Phi_3(q)=q^2+q+1$ & $\Phi_3$            & $13$ \\
$\Phi_4(q)=q^2+1$   & $\Phi_4$            & $10$ \\
$\Phi_6(q)=q^2-q+1$ & $\Phi_6$            & $7$ \\
$\Sigma=2^{\Phi_4/2}$& $\Sigma$            & $32$ \\
$P=q(\Phi_4/2)(\Phi_4{+}\Phi_6)$ & $P$ & $255$ \\
\bottomrule
\end{tabular}
\caption{Parameters of $W(3,3)$ at $q=3$.}
\end{table}

The Ihara zeta function is
\[
  \zeta_{W}(u)^{-1} = (1-u^2)^{k(v-1)/2}\,p_1(u)^f\,p_2(u)^g,
\]
with $p_1(u)=1-2u+11u^2$, $p_2(u)=1+4u+11u^2$,
and discriminants $\mathrm{disc}(p_1)=-4\Phi_4=-40$,
$\mathrm{disc}(p_2)=-4\Phi_6=-28$.

%%% SECTION 2 %%%
\section{The W(3,3) Uniqueness Theorem}
\label{sec:uniqueness}

Let $q$ be a prime power and $W(3,q)$ the Lubotzky--Phillips--Sarnak
Ramanujan bipartite graph with degree $k=q(q+1)$,
vertex count $v=(q+1)(q^2+1)$ per side, and adjacency eigenvalues
$r = q-1$ and $s = -(q+1)$, with multiplicities
$f = q(q+1)^2/2$ and $g = q(q^2+1)/2$.
Define
\[
  \Phi_3(q) = q^2+q+1,\qquad
  \Phi_4(q) = q^2+1,\qquad
  \Phi_6(q) = q^2-q+1.
\]
The Ihara zeta function factors through the two quadratic polynomials
\[
  p_1(u) = 1 - (q-1)u + (k-1)u^2,\qquad
  p_2(u) = 1 + (q+1)u + (k-1)u^2,
\]
whose shared constant term is $k(q)-1 = q^2+q-1$.

\begin{definition}[Spectral deficit]
Define the spectral deficit by
\[
  \mathrm{Deficit}(q) = -\frac{\mathrm{disc}(p_1)}{4} - \Phi_4(q).
\]
Since
\[
  \mathrm{disc}(p_1) = (q-1)^2 - 4(k-1) = -3q^2-6q+5,
\]
we obtain the exact factorization below.
\end{definition}

\begin{lemma}[C1]
\[
  \mathrm{Deficit}(q) = -\frac{(q-3)^2}{4}.
\]
Hence $\mathrm{Deficit}(q)=0$ if and only if $q=3$.
\end{lemma}
\begin{proof}
A direct expansion gives
\[
  \mathrm{Deficit}(q)
  = \frac{3q^2+6q-5}{4} - (q^2+1)
  = \frac{-q^2+6q-9}{4}
  = -\frac{(q-3)^2}{4}. \qedhere
\]
\end{proof}

\begin{lemma}[C2]
$k(q)+g(q)=q^q$ if and only if $q=3$.
\end{lemma}
\begin{proof}
We compute
\[
  k(q)+g(q)=q(q+1)+\frac{q(q^2+1)}{2} = \frac{q^3}{2}+q^2+\frac{3q}{2}.
\]
At $q=3$ this equals $27=3^3$.
For $q\ge 4$, the right-hand side $q^q$ dominates any cubic polynomial,
so no further positive integer solutions exist.
\end{proof}

\begin{lemma}[C3]
The identities
\[
  -f(q)=\tau(2),\qquad k(q)\,q\,\Phi_6(q)=\tau(3)
\]
hold if and only if $q=3$.
\end{lemma}
\begin{proof}
Since $\tau(2)=-24$ and $\tau(3)=252$, we evaluate at $q=3$:
\[
  -f(3)=-\frac{3\cdot 4^2}{2}=-24,
  \qquad
  k(3)\cdot 3\cdot \Phi_6(3)=12\cdot 3\cdot 7=252.
\]
The first equation has unique positive integer solution $q=3$,
and the second is strictly increasing for $q\ge 2$.
\end{proof}

\begin{lemma}[C4]
$\mathbb{Q}(\sqrt{-\Phi_6(q)})$ has class number $1$ for $q\in\{2,3,7\}$.
In conjunction with C1, this is unique to $q=3$.
\end{lemma}
\begin{proof}
The Heegner discriminants are
\(
\{1,2,3,4,7,8,11,19,43,67,163\}.
\)
Now
\[
  \Phi_6(2)=3,\qquad \Phi_6(3)=7,\qquad \Phi_6(7)=43,
\]
all of which are Heegner. But C1 fails at $q=2$ and $q=7$.
\end{proof}

\begin{lemma}[C5]
Let $E_{-\Phi_6(q)}$ be the CM elliptic curve attached to the order of
 discriminant $-\Phi_6(q)$. Then
\[
  a_{k(q)-1}(E_{-\Phi_6(q)}) = \mathrm{ev}_s(q)=-(q+1)
\]
holds at $q=3$ and fails already at $q=2$.
\end{lemma}
\begin{proof}
At $q=3$, one has $k-1=11$, $\Phi_6=7$, and
\[
  a_{11}(E_{-7})=-4=-(3+1)=\mathrm{ev}_s(3).
\]
At $q=2$, one has $k-1=5$, $\Phi_6=3$, while direct point-counting gives
\[
  a_5(E_{-3})=0\neq -3=\mathrm{ev}_s(2).
\]
So the identity is not generic in the $W(3,q)$ family.
\end{proof}

\begin{theorem}[W(3,3) Uniqueness]
$W(3,3)$ is the unique member of the $W(3,q)$ family satisfying
all five characterizations C1--C5 simultaneously.
\end{theorem}
\begin{proof}
By the preceding lemmas, C1, C2, and C3 each isolate $q=3$,
while C1$\wedge$C4 and C5 also hold uniquely at $q=3$.
Therefore the full conjunction C1$\wedge$C2$\wedge$C3$\wedge$C4$\wedge$C5
holds if and only if $q=3$.
\end{proof}


%%% SECTION 3 %%%
\section{The Ramanujan Tau Bridge}
\label{sec:tau}

Let $\tau\colon\mathbb{Z}_{>0}\to\mathbb{Z}$ be the Ramanujan tau function,
defined by
\[
  \Delta(q) = q\prod_{n=1}^{\infty}(1-q^n)^{24} = \sum_{n=1}^{\infty}\tau(n)q^n,
  \qquad \tau(2)=-24,\;\tau(3)=252.
\]

\begin{theorem}[Tau bridge]\label{thm:tau}
$\tau(2) = -f(3) = -24$ and $\tau(3) = k(3)\cdot 3\cdot\Phi_6(3) = 252$.
\end{theorem}
\begin{proof}
$-f(3)=-q(q+1)^2/2\big|_{q=3}=-24$. And $k\cdot q\cdot\Phi_6=12\cdot 3\cdot 7=252$.
\end{proof}

\begin{remark}[Hecke recursion]
The recursion $\tau(p^{a+1})=\tau(p)\tau(p^a)-p^{11}\tau(p^{a-1})$
at $p=k-1=11$ involves $p^{11}=(k-1)^{11}$, connecting
the weight-$12=k$ structure of $\Delta$ to the Ramanujan degree $k-1=11$.
\end{remark}

\begin{theorem}[$\Phi_6$ barrier]\label{thm:barrier}
For primes $p\equiv 1\pmod{\Phi_6(3)=7}$, Ramanujan's congruence
$\tau(p)\equiv 0\pmod{p}$ holds.
\end{theorem}

\begin{theorem}[C5 via Frobenius]\label{thm:C5}
Let $E\colon y^2=x^3-35x+98$ \emph{(}$j=-3375=-g^3$, conductor $49$\,\emph{)}.
This curve has CM by $\mathbb{Z}[(1+\sqrt{-7})/2]$~\cite{LMFDB};
for inert primes $\left(\frac{-7}{p}\right)=-1$ one has $a_p(E)=0$, and
\[
  a_{k-1}(E) = a_{11}(E) = -(q+1) = \mathrm{ev}_s.
\]
The discriminant identity $a_{11}^2-4\cdot 11=16-44=-28=-4\Phi_6$
identifies $p_2(u)=1+\mathrm{ev}_s\cdot u+(k-1)u^2$ as the
Euler factor of $L(E,s)$ at $p=k-1$.
\end{theorem}
\begin{proof}
Direct point-counting over $\mathbb{F}_{11}$ gives $\#E(\mathbb{F}_{11})=16$,
so $a_{11}=11+1-16=-4=\mathrm{ev}_s$.
The CM property (quadratic twist of $y^2=x^3-35x-98$ by $d=-1$~\cite{LMFDB})
ensures $a_p=0$ for all inert $p$, verified for $p\in\{3,5,13,17,19\}$.
\end{proof}

\subsection*{The $j$-tower and its cube structure}

\begin{remark}[C6: divisibility of $g$ by $3$]\label{rem:C6}
Among all prime powers $q\equiv 1\pmod 4$, $3\mid g(q)=q(q^2+1)/2$
if and only if $3\mid q$; for prime $q$ this forces $q=3$ uniquely.
\begin{proof}
$3\mid g(q)\Leftrightarrow 3\mid q(q^2+1)$.
If $q\equiv 1\pmod 3$: $q(q^2+1)\equiv 2$;
if $q\equiv 2\pmod 3$: $q(q^2+1)\equiv 1$; neither is $\equiv 0$.
\end{proof}
Set $\Sigma=2^{g/3}=2^{\Phi_4/2}=2^5=32$ (well-defined since $3\mid g=15$
and $g/3=\Phi_4/2=5$). The identity $g/3=\Phi_4/2$ is a cross-sector link:
the $D=-11$ cube root is determined by the $r$-sector parameter $\Phi_4$.
\end{remark}

The Ihara zeta function factors through
$\mathrm{disc}(p_2)=-4\Phi_6=-28$ ($s$-sector, class number $1$, Heegner) and
$\mathrm{disc}(p_1)=-4\Phi_4=-40$ ($r$-sector, class number $2$).
Let $P=q(\Phi_4/2)(\Phi_4+\Phi_6)=3\cdot 5\cdot 17=255$.
The complete $j$-tower~\cite{LMFDB} is:

\begin{equation}\label{eq:jtower}
\renewcommand{\arraystretch}{1.3}
\begin{array}{rcllll}
  j(D=-4)  &=& k^3 = 1728         &k=q(q+1)         &\mathbb{Z}[i]\\
  j(D=-7)  &=& -g^3 = -3375       &g=q(q^2+1)/2     &\mathbb{Z}[\tfrac{1+\sqrt{-7}}{2}]\\
  j(D=-8)  &=& (v/2)^3 = 8000     &v/2=(q+1)(q^2+1)/2&\mathbb{Z}[\sqrt{-2}]\\
  j(D=-11) &=& -\Sigma^3 = -32768 &\Sigma=2^{\Phi_4/2}&\mathbb{Z}[\tfrac{1+\sqrt{-11}}{2}]\\
  j(D=-28) &=& P^3=16{,}581{,}375 &P=q\tfrac{\Phi_4}{2}(\Phi_4+\Phi_6)&\mathbb{Z}[\sqrt{-7}],f=2
\end{array}
\end{equation}

\begin{proposition}[Cube structure]\label{prop:cube}
Every $j$-tower entry satisfies $j(D)=\pm(\text{spectral invariant})^3$.
No Heegner discriminant with $|D|>4\Phi_6=28$ has its cube root
in $\{k,\,g,\,v/2,\,\Sigma,\,P\}$.
Equivalently, the tower is the complete intersection of:
\textnormal{(a)} Heegner discriminants with $|D|\leq 4\Phi_6$;\;
\textnormal{(b)} the cube root is a $W(3,3)$ spectral invariant.
\end{proposition}

\noindent
The $r$-eigenspace sector ($h=2$, disc $-40$) has $j$-values
$\approx 4.3\times 10^8$ and $\approx 2.1\times 10^4$ that are not
cubes of any $W(3,3)$ parameter, confirming the $j$-tower is entirely
an $s$-sector phenomenon.


%%% SECTION 4 %%%
\section{Neutrino Mass Constraint}
\label{sec:neutrino}

For a neutrino mass vector $\mathbf{m}=(m_1,m_2,m_3)$ normalized to $m_{\max}=1$,
define
\[
  \mu_{\rm eff}^2(\mathbf{m}) = -\frac{\ln(\prod_i m_i/m_{\max})^{1/3}}{\ln\Phi_4}.
\]

\begin{theorem}[Entry point]
The $W(3,3)$ spectral constraint $\mu_{\rm eff}^2=1/(q+1)=1/4$
has a unique solution $m_1=m_1^*$ for each neutrino mass ordering.
Under the normal hierarchy with NuFIT 5.3 oscillation parameters \cite{NuFIT53},
$m_1^*\approx 50.7\,\mathrm{meV}$ and $\sum m_\nu\approx 167\,\mathrm{meV}$.
\end{theorem}

\begin{corollary}[DESI consistency]
The DESI DR1 bound $\sum m_\nu < 72\,\mathrm{meV}$ \cite{DESI2024}
excludes the NH entry point at greater than $2\sigma$.
The $W(3,3)$ spectral constraint predicts quasi-degenerate neutrinos
in tension with current cosmological data.
\end{corollary}

%%% SECTION 5 %%%
\section{The Seesaw Cascade}
\label{sec:seesaw}

Define the seesaw map $T\colon\mu\mapsto\mu_{\rm eff}^2(1/\mathbf{m}(\mu))$.

\begin{theorem}[Cascade ratio]\label{thm:cascade}
In the quasi-degenerate NH limit,
\[
  r = \frac{\Delta m^2_{21}+\Delta m^2_{31}}{2\Delta m^2_{31}-\Delta m^2_{21}}
  \;\xrightarrow{\Delta m^2_{21}\to 0}\;\frac{1}{2},
  \qquad r-\tfrac{1}{2} = \frac{3\Delta m^2_{21}}{2(2\Delta m^2_{31}-\Delta m^2_{21})}.
\]
With NuFIT 5.3 values~\cite{NuFIT53}: $r=0.52244$.
\end{theorem}
\begin{proof}
Expand $\mu_{\rm eff}^2$ to leading order in $\Delta m^2/m_1^2$;
the ratio $r=B/A$ with $A=(\Delta m^2_{31}/3-\Delta m^2_{21}/6)/\ln\Phi_4$
and $B=(\Delta m^2_{21}+\Delta m^2_{31})/(6\ln\Phi_4)$.
$\Phi_4$ cancels; $r-1/2=3A'/(2(2B'-A'))$ follows algebraically.
\end{proof}

The first cascade iterates land near $W(3,3)$ values:
$T^0=1/4=1/m$, $T^1\approx 1/\Phi_6$, $T^2\approx 1/\Phi_3$;
asymptotically $T^{n+1}/T^n\to r$.

\begin{thebibliography}{99}
\bibitem{LPS1988}
  A.~Lubotzky, R.~Phillips, P.~Sarnak,
  \emph{Ramanujan graphs},
  Combinatorica \textbf{8} (1988), 261--277.

\bibitem{Deligne1974}
  P.~Deligne,
  \emph{La conjecture de Weil I},
  Publ.\ Math.\ IHES \textbf{43} (1974), 273--307.

\bibitem{LMFDB}
  The LMFDB Collaboration,
  \emph{The $L$-functions and Modular Forms Database},
  \texttt{https://www.lmfdb.org}, 2024.

\bibitem{NuFIT53}
  I.~Esteban, M.~C.~Gonzalez-Garcia, M.~Maltoni, T.~Schwetz, A.~Zhou,
  JHEP \textbf{09} (2020) 178; \texttt{http://www.nu-fit.org}.

\bibitem{DESI2024}
  DESI Collaboration,
  \emph{DESI 2024 VI: Cosmological constraints from BAO},
  arXiv:2404.03002.

\bibitem{Dahn2026repo}
  W.~Dahn,
  \emph{$W(3,3)$-Theory: Automated tests and computations},
  \texttt{https://github.com/wilcompute/W33-Theory}, 2026.
\end{thebibliography}

\end{document}
